\documentclass[12pt]{article}
\usepackage[utf8]{inputenc}
\usepackage{amsmath, amssymb, geometry, graphicx, setspace}
\geometry{letterpaper, margin=1in}
\setstretch{1.15}

\title{\textbf{Interplanetary Trajectories: \\ A Step-by-Step Guide to the Patched Conic Method}}
\author{Module: Space Engineering Vehicle and Mission Design}
\date{}

\begin{document}

\maketitle

\section*{The Challenge of Interplanetary Flight}
An interplanetary mission is fundamentally an $N$-body problem. A spacecraft traveling from Earth to Mars is gravitationally influenced simultaneously by the Earth, the Sun, Mars, and other celestial bodies. Because there is no general analytical solution to the $N$-body problem, we rely on a powerful simplification known as the \textbf{Patched Conic Approximation}.

This method relies on the fact that gravity obeys an inverse-square law. Close to a planet, the planet's gravity completely dominates the spacecraft's motion. Far from a planet, the Sun's gravity dominates. By dividing the solar system into distinct gravitational zones, we break the complex journey into three solvable $2$-body problems (conic sections) and ``patch'' them together at their boundaries.

\vspace{0.5cm}

\subsection*{The Sphere of Influence (SOI)}
The boundary where we switch our primary frame of reference from the planet to the Sun (or vice versa) is called the \textit{Sphere of Influence} ($r_{SOI}$). At this distance, the perturbing effect of the Sun on the planet-centric orbit roughly equals the perturbing effect of the planet on the Sun-centric orbit. 
\begin{equation}
    r_{SOI} = R_p \left( \frac{m_{planet}}{m_{sun}} \right)^{2/5}
\end{equation}
where $R_p$ is the distance from the Sun to the planet. While the SOI is vast from a human perspective, it is a mere speck on the scale of the solar system. Therefore, when looking at the Sun-centric leg of the journey, the SOI is treated simply as a point coinciding with the planet's center.

\newpage

% ---------------------------------------------------------
% THEORETICAL STEPS
% ---------------------------------------------------------
\section*{Step 1: The Heliocentric Transfer (Phase 2)}
In this phase, we ignore the planets' gravity and treat them as moving points. The spacecraft travels on a heliocentric transfer ellipse from the orbit of Planet 1 ($R_1$) to the orbit of Planet 2 ($R_2$). 

For a Hohmann transfer, the semi-major axis is $a_{tx} = (R_1 + R_2)/2$. The absolute velocity the spacecraft must have with respect to the Sun ($V_{tx}$) at departure is:
\begin{equation}
    V_{tx} = \sqrt{\mu_{sun} \left( \frac{2}{R_1} - \frac{1}{a_{tx}} \right)}
\end{equation}

Planet 1 is moving around the Sun with velocity $V_1 = \sqrt{\mu_{sun}/R_1}$. The \textbf{Hyperbolic Excess Velocity ($v_\infty$)} is the relative speed the spacecraft must possess as it crosses the boundary of Planet 1's SOI:
\begin{equation}
    v_{\infty} = | V_{tx} - V_1 |
\end{equation}

\section*{Step 2: Planetocentric Departure (Phase 1)}
Assuming the spacecraft starts in a circular parking orbit of radius $r_p$, its initial velocity is $v_c = \sqrt{\mu_1 / r_p}$. By conservation of energy, the required velocity at the periapsis of the departure hyperbola is:
\begin{equation}
    v_p = \sqrt{v_{\infty}^2 + \frac{2\mu_1}{r_p}}
\end{equation}

The engine burn required is:
\begin{equation}
    \Delta V_{depart} = v_p - v_c = \sqrt{v_{\infty}^2 + \frac{2\mu_1}{r_p}} - \sqrt{\frac{\mu_1}{r_p}}
\end{equation}

\section*{Step 3: Planetocentric Arrival (Phase 3)}
At Planet 2, the arrival velocity relative to the Sun is $V_{arr}$. Planet 2 is moving with velocity $V_2 = \sqrt{\mu_{sun}/R_2}$. The hyperbolic excess speed at Planet 2 is $v_{\infty, arr} = | V_{arr} - V_2 |$.

To capture into a circular orbit of radius $r_{cap}$, we calculate the speed at periapsis of the arrival hyperbola and subtract the target circular speed:
\begin{equation}
    v_{p, arr} = \sqrt{v_{\infty, arr}^2 + \frac{2\mu_2}{r_{cap}}} \quad , \quad v_{c, arr} = \sqrt{\frac{\mu_2}{r_{cap}}}
\end{equation}
\begin{equation}
    \Delta V_{capture} = v_{p, arr} - v_{c, arr}
\end{equation}

\newpage

% ---------------------------------------------------------
% WORKED EXAMPLE
% ---------------------------------------------------------
\section*{Worked Numerical Example: Earth to Mars}
Let us design a trajectory from a 300 km altitude Earth parking orbit to a 300 km altitude Mars capture orbit. We assume circular, coplanar planetary orbits.

\textbf{Given Constants:}
\begin{itemize}
    \item $\mu_{sun} = 1.327 \times 10^{11} \text{ km}^3/\text{s}^2$
    \item $R_{\oplus} = 1.496 \times 10^8 \text{ km}$, $\quad \mu_{\oplus} = 3.986 \times 10^5 \text{ km}^3/\text{s}^2$, $\quad r_{p} = 6678 \text{ km}$
    \item $R_{Mars} = 2.279 \times 10^8 \text{ km}$, $\quad \mu_{Mars} = 4.282 \times 10^4 \text{ km}^3/\text{s}^2$, $\quad r_{cap} = 3696 \text{ km}$
\end{itemize}

\subsection*{1. Heliocentric Transfer}
The transfer semi-major axis $a_{tx} = (1.496 + 2.279) \times 10^8 / 2 = 1.8875 \times 10^8 \text{ km}$.

Earth's velocity: $V_{\oplus} = \sqrt{\frac{1.327 \times 10^{11}}{1.496 \times 10^8}} = 29.78 \text{ km/s}$.

Required transfer velocity at Earth: 
\begin{equation*}
    V_{tx} = \sqrt{1.327 \times 10^{11} \left( \frac{2}{1.496 \times 10^8} - \frac{1}{1.8875 \times 10^8} \right)} = 32.73 \text{ km/s}
\end{equation*}

Earth departure hyperbolic excess velocity: $v_{\infty} = 32.73 - 29.78 = \textbf{2.95 km/s}$.

\subsection*{2. Earth Departure Burn}
Parking orbit velocity: $v_c = \sqrt{\frac{3.986 \times 10^5}{6678}} = 7.73 \text{ km/s}$.

Required burnout velocity at periapsis:
\begin{equation*}
    v_p = \sqrt{2.95^2 + \frac{2(3.986 \times 10^5)}{6678}} = \sqrt{8.70 + 119.38} = 11.32 \text{ km/s}
\end{equation*}

Departure burn: $\Delta V_{depart} = 11.32 - 7.73 = \textbf{3.59 km/s}$.

\subsection*{3. Mars Arrival Burn}
Mars velocity: $V_{Mars} = \sqrt{\frac{1.327 \times 10^{11}}{2.279 \times 10^8}} = 24.13 \text{ km/s}$.

Arrival transfer velocity at Mars:
\begin{equation*}
    V_{arr} = \sqrt{1.327 \times 10^{11} \left( \frac{2}{2.279 \times 10^8} - \frac{1}{1.8875 \times 10^8} \right)} = 21.48 \text{ km/s}
\end{equation*}

Mars arrival hyperbolic excess velocity: $v_{\infty, arr} = 24.13 - 21.48 = \textbf{2.65 km/s}$.

Capture orbit velocity: $v_{c, arr} = \sqrt{\frac{4.282 \times 10^4}{3696}} = 3.40 \text{ km/s}$.

Velocity at periapsis of arrival hyperbola:
\begin{equation*}
    v_{p, arr} = \sqrt{2.65^2 + \frac{2(4.282 \times 10^4)}{3696}} = \sqrt{7.02 + 23.17} = 5.49 \text{ km/s}
\end{equation*}

Capture burn: $\Delta V_{capture} = 5.49 - 3.40 = \textbf{2.09 km/s}$.

\subsection*{Total Mission $\Delta V$}
\begin{equation*}
    \Delta V_{total} = 3.59 \text{ km/s} + 2.09 \text{ km/s} = \textbf{5.68 km/s}
\end{equation*}

\end{document}
